\section{Repeated stable summands}\label{sec:repeated}

The relative estimate also applies when some stable summands are
isomorphic. The limiting kernel then contains matrices acting between
the repeated copies, and the comparison operator is a commutator
quadratic form on this larger space.

Retain the curve, scaling exponents, and harmonic upper-triangular family
\eqref{eq:several-family}, but allow repetitions among the stable
equal-slope bundles $F_1,\ldots,F_m$. Choose pairwise nonisomorphic
representatives $H_1,\ldots,H_s$ and identifications
\begin{equation}\label{eq:isotypical-splitting}
 G=\bigoplus_{i=1}^m F_i
   \simeq\bigoplus_{\alpha=1}^s H_\alpha\otimes\CC^{m_\alpha},
 \qquad \sum_\alpha m_\alpha=m.
\end{equation}
Here $H_\alpha\otimes\CC^{m_\alpha}$ is the direct sum of all copies
of $H_\alpha$, called its isotypical component. Fix a
Hermitian--Einstein metric on each $H_\alpha$ and give its copies the
transported metric. Thus the product metric on each isotypical
component is $h_{H_\alpha}\otimes I$. Write
$r_\alpha=\rk H_\alpha$.

Stability at equal slope gives
\begin{equation}\label{eq:matrix-algebra}
 \mathcal A:=H^0(X,\End G)
 =\bigoplus_{\alpha=1}^s
   \id_{H_\alpha}\otimes M_{m_\alpha}(\CC).
\end{equation}
Every element of $\mathcal A$ is parallel for the split Chern
connection: under the chosen identifications its nonzero entries
are constant multiples of the identity of a stable representative.
Consequently
\begin{equation}\label{eq:matrix-kernel}
 \begin{split}
 \mathcal K:=H^0(X,\End_0G)
 &=\left\{(X_\alpha)_\alpha:
             \sum_\alpha r_\alpha\tr X_\alpha=0\right\},\\
 d:=\dim_\CC\mathcal K&=\sum_\alpha m_\alpha^2-1,\qquad
 \lVert X\rVert^2=V\sum_\alpha r_\alpha\tr(X_\alpha^*X_\alpha).
 \end{split}
\end{equation}
Use $D_t=D_0+B_t$, $B_tx=[N_t,x]$, $\Delta_t=D_t^*D_t$,
$P,Q,\gamma,\eta_t$, and $p_*$ as before, now with this kernel.
Define the nonnegative self-adjoint operator $L_t$ on $\mathcal K$ by
\begin{equation}\label{eq:matrix-form}
 \langle x,L_tx\rangle=A_t(x):=\lVert[N_t,x]\rVert_{L^2}^2,
 \qquad L_t=PB_t^*B_tP|_{\mathcal K}.
\end{equation}
Let $\nu_1(t)\le\cdots\le\nu_d(t)$ denote its eigenvalues, including
zeros. Unlike \eqref{eq:graph-form}, this form may have cross terms
between extension entries contributing to the same matrix block.

\begin{theorem}[Comparison with the commutator form]\label{thm:matrix-relative}
If $\eta_t\le\sqrt\gamma/4$ and
$\gamma_t=(\sqrt\gamma-\eta_t)^2$, then exactly $d$ eigenvalues of
$\Delta_t$ lie below $\gamma/4$, counting zeros, and
\begin{equation}\label{eq:matrix-relative}
 \left(1-\frac{2\eta_t^2}{\gamma_t}\right)\nu_j(t)
 \le\lambda_j(t)\le\nu_j(t),\qquad 1\le j\le d,
 \qquad \lambda_{d+1}(t)\ge\frac\gamma4.
\end{equation}
In particular, $\lambda_j(t)/\nu_j(t)=1+O(t^{2p_*})$ for every
positive $\nu_j(t)$, uniformly over these $d$ modes.

For every $t>0$, the trace-free holomorphic endomorphism space of
$E_t$ has dimension
\begin{equation}\label{eq:matrix-centralizer}
 \dim_\CC\{x\in\mathcal K:[N_t,x]=0\}.
\end{equation}
Thus $E_t$ is simple if and only if the only elements of $\mathcal A$
commuting with $N_t$ are scalar identities.

For a unit small eigensection with eigenvalue $\lambda$, and for the
orthogonal projection $\Pi_t$ onto the full $d$-dimensional small
spectral subspace, one has, for every $k\ge0$,
\begin{equation}\label{eq:matrix-convergence}
 \lVert Qu\rVert_{C^k}\le C_k t^{p_*}\sqrt\lambda,
 \qquad
 \lVert\Pi_t-P\rVert_{L^2\to C^k}=O(t^{2p_*}).
\end{equation}
\end{theorem}

\begin{proof}
Multiplication on either side by a parallel endomorphism commutes
with the split Dolbeault Laplacian on bundle-valued forms. Since
$N_t$ is harmonic and $x\in\mathcal K$ is parallel, $B_tx=[N_t,x]$
is harmonic. In particular,
\[
 D_0^*B_tx=0\qquad(x\in\mathcal K).
\]
For $u=x+y$, $x\in\mathcal K$ and $y\in QH^1$, the same gap
estimate and cross-term bound as in \eqref{eq:complement-gap} give
\[
 \lVert D_tu\rVert^2
 \ge\left(1-\frac{2\eta_t^2}{\gamma_t}\right)A_t(x)
       +\frac{\gamma_t}{2}\lVert y\rVert^2.
\]
Also $\nu_d(t)\le\eta_t^2\le\gamma/16$. The min--max argument
in Theorem~\ref{thm:relative} therefore applies with $d$ in place
of $m-1$ and proves \eqref{eq:matrix-relative}. No commutativity of
$\mathcal A$ is used in this argument.

The displayed lower bound identifies the kernel for small $t$ with
$\{x\in\mathcal K:B_tx=0\}$. The parallel diagonal map $g_t$
preserves $\mathcal A$ by conjugation, preserves trace, and satisfies
$N_t=g_tN_1g_t^{-1}$. It therefore identifies the commutants for
different positive parameters. Since the bundles $E_t$ are also
holomorphically isomorphic, the kernel dimension formula holds for
all $t>0$.

To check smooth convergence, the complementary equation is still
\eqref{eq:schur}, with
$R_t=QB_t^*B_tP$ and $\lVert R_tx\rVert\le\eta_t\sqrt{A_t(x)}$.
Every $B_tx$ lies in the fixed finite-dimensional space of harmonic
$(0,1)$-forms with values in $\End_0G$. Equivalence of norms on
that space gives
\[
 \lVert B_tx\rVert_{H^\ell}\le C_\ell\sqrt{A_t(x)},\qquad
 \lVert R_tx\rVert_{H^\ell}
       \le C_\ell t^{p_*}\sqrt{A_t(x)}.
\]
This estimate remains valid in the presence of cancellations between
matrix entries. The uniform complementary inverse estimates and
Sobolev embedding used in Proposition~\ref{prop:subspace} now give
the first assertion of \eqref{eq:matrix-convergence}. The largest
small eigenvalue is $O(t^{2p_*})$; applying that assertion to an
orthonormal basis of the $d$ small eigensections and using the
projection formula for the resulting graph over $\mathcal K$
proves the second assertion.
\end{proof}

When all $m_\alpha=1$, \eqref{eq:matrix-form} is exactly the
weighted graph form. For repeated factors, connectedness of that
graph is insufficient for simplicity. For example, take
$G=H\oplus H$ with $H$ stable and
$N_t=t\beta E_{12}$, where $\beta$ is a nonzero harmonic scalar
$(0,1)$-form. The parallel nilpotent endomorphism $E_{12}$ commutes
with $N_t$, so this connected two-vertex family is not simple.
Here $E_{ij}$ denotes the identity map from the $j$th copy of $H$
to the $i$th, extended by zero on the other blocks.

The mean-curvature calculation \eqref{eq:several-curvature} uses
only harmonicity and the common Einstein constant, so it also applies
with repeated factors. In particular, the transported metrics
$g_t^*h_0$ have uniform mean-curvature error $O(t^{2p_*})$ on the
fixed bundle $E_1$.

\subsection{Recovery of factors and multiplicities}

After adjoining the scalar identity, the limiting small subspace is
the algebra $\mathcal A$ in \eqref{eq:matrix-algebra}. Its minimal
nonzero central idempotents are the projections onto the isotypical
components. The corresponding ideals have dimensions $m_\alpha^2$,
so their matrix sizes determine the multiplicities. Within such an
ideal, a minimal nonzero idempotent has rank one on the multiplicity
space and its holomorphic image is isomorphic to $H_\alpha$.
A decomposition of the central identity into orthogonal rank-one
projections gives $m_\alpha$ copies. Such a decomposition is not
canonical when $m_\alpha>1$.

The abstract algebra alone determines the multiplicities, but not the
isomorphism classes of the bundles $H_\alpha$; these come from its action
on $G$.

\subsection{A simple example with a repeated line bundle}\label{sec:repeated-example}

Let $X$ have genus at least two and volume one. Choose pairwise
nonisomorphic degree-zero line bundles $A,B,C$, with their
Hermitian--Einstein metrics, and put
$G=A\oplus B\oplus A\oplus C$, using the same metric on both
copies of $A$. Riemann--Roch gives
$h^1(X,L)=g(X)-1>0$ for every nontrivial degree-zero line bundle
$L$. Thus one can choose harmonic forms of unit $L^2$ norm
\[
 a\in\Omega^{0,1}(\Hom(B,A)),\quad
 b\in\Omega^{0,1}(\Hom(A,B)),\quad
 c\in\Omega^{0,1}(\Hom(C,A)).
\]
Set $N_{12}=ta$, $N_{23}=t^2b$, $N_{34}=tc$, and all other entries
to zero. These scales come from the scaling exponents $(4,3,1,0)$. An element of
$\mathcal K$ is
\[
 x=\begin{pmatrix}
 x_1&0&u&0\\0&x_2&0&0\\v&0&x_3&0\\0&0&0&x_4
 \end{pmatrix},\qquad \sum_{i=1}^4x_i=0,
 \qquad \lVert x\rVert^2=\sum_i|x_i|^2+|u|^2+|v|^2.
\]
The six nonzero commutator entries give
\begin{equation}\label{eq:repeated-example-form}
 \begin{split}
 A_t(x)={}&t^2\bigl(|x_2-x_1|^2+|x_4-x_3|^2+|u|^2+|v|^2\bigr)\\
         &+t^4\bigl(|x_3-x_2|^2+|v|^2\bigr).
 \end{split}
\end{equation}
This is positive definite on $\mathcal K$ for $t>0$, so
Theorem~\ref{thm:matrix-relative} proves that $E_t$ is simple.
Its upper-triangular filtration has stable degree-zero quotients
and a degree-zero line subbundle; hence it is strictly semistable.

The $u$ and $v$ modes have comparison eigenvalues $t^2$ and
$t^2+t^4$. The diagonal part is the four-vertex path with
conductances $t^2,t^4,t^2$. Its three positive eigenvalues are
\[
 2t^2,\qquad t^2+t^4\pm\sqrt{t^4+t^8}.
\]
Thus there are $2^2+1^2+1^2-1=5$ small positive Dolbeault
eigenvalues, in increasing order,
\begin{equation}\label{eq:repeated-example-rates}
 \begin{aligned}
 \lambda_1&=t^4(1+O(t^2)),\\
 \lambda_j&=t^2(1+O(t^2))&& (j=2,3),\\
 \lambda_j&=2t^2(1+O(t^2))&& (j=4,5).
 \end{aligned}
\end{equation}
The full limiting space contains the two off-diagonal maps between
the copies of $A$, in addition to the three trace-free diagonal
directions. After adjoining the identity, it is
$M_2(\CC)\oplus\CC\oplus\CC$, realized on $G$.

For completeness the limiting subspaces at each leading coefficient
can also be read from this example. The slow line is spanned by
$\diag(1,1,-1,-1)$. The coefficient-one subspace at order $t^2$
is spanned by $E_{13},E_{31}$, and the coefficient-two subspace
is spanned by $\diag(1,-1,0,0)$ and $\diag(0,0,1,-1)$.
Indeed, $t^{-2}L_t$ converges to the operator with exactly these
eigenspaces and respective eigenvalues $0,1,2$. In the complementary
equation \eqref{eq:schur}, the correction has operator norm
$O(t^4)$ by \eqref{eq:schur-bound}, which applies equally to
\eqref{eq:matrix-form}. After division by $t^2$ it tends to zero.
Together with \eqref{eq:matrix-convergence} and the known
multiplicities, this proves convergence of the three spectral projections
in every $L^2\to C^k$ norm. It does not single out a basis within
either two-dimensional limiting subspace.
